\section{Introduction}
\label{sec:intro}

Parton physics is important both for understanding the dynamics of
high-energy collisions of hadrons and for studying their
internal structure~\cite{Ellis:1991qj,Thomas:2001kw}. The most familiar examples
are quark and gluon parton distribution functions (PDFs)
which, on one hand, provide the beam information for
high-energy productions at colliders~\cite{Gao:2017yyd}, and on the other hand,
describe the bound-state physics of the colliding hadrons.

Despite its importance, calculating parton physics from first principles of
quantum chromodynamics (QCD) has been a challenge. Recently,
an effective field theory (EFT) approach--large momentum effective theory (LaMET)--has been proposed to extract parton physics from physical properties
of hadrons moving at large momentum~\cite{Ji:2013dva,Ji:2014gla}, where the latter
can be calculated from systematic approximations to Euclidean QCD such
as lattice field theory. Since its proposal, LaMET has been widely used in calculating quark isovector distribution functions~\cite{Lin:2014zya,Alexandrou:2015rja,Chen:2016utp,Alexandrou:2016jqi,Alexandrou:2018pbm,Chen:2018xof,Lin:2018qky,Liu:2018uuj,Alexandrou:2018eet,Liu:2018hxv,Chen:2018fwa,Izubuchi:2019lyk,Shugert:2020tgq,Chai:2020nxw,Lin:2020ssv,Fan:2020nzz}, distribution amplitudes (DAs)~\cite{Zhang:2017bzy,Chen:2017gck,Zhang:2020gaj},
generalized parton distributions~\cite{Chen:2019lcm,Alexandrou:2019dax}, and recently transverse-momentum-dependent
distributions~\cite{Shanahan:2019zcq,Shanahan:2020zxr,Zhang:2020dbb}, and even higher-twist distributions~\cite{Bhattacharya:2020cen}. Some recent reviews on
LaMET can be found in Refs.~\cite{Ji:2020ect,Cichy:2018mum}.

The key idea of LaMET is that partons in the infinite momentum frame (IMF)
can be approximated by physical properties of a hadron at large but
finite momentum. Due to the existence of ultraviolet (UV) divergences, this approximation
is not completely straightforward. It requires using the standard EFT technology of matching and running. Detailed investigations have shown that the standard DGLAP evolution~\cite{Dokshitzer:1977sg,Gribov:1972ri,Altarelli:1977zs} has its origin in the momentum evolution of physical properties of the hadron~\cite{Ji:2020ect}.

In LaMET applications, one begins with lattice calculations of spatial correlation functions. Since lattice
breaks the continuum symmetry, power divergences
appear in bare correlation functions. They must be subtracted
when matched to those in a continuum scheme such as dimensional regularization
and (modified)-minimal subtraction, ${\overline {\rm MS}}$. In the past, the main
approaches suggested in practical applications include the regularization-independent momentum subtraction method
or RI/MOM~\cite{Constantinou:2017sej,Stewart:2017tvs,Alexandrou:2017huk,Chen:2017mzz} and the
ratio method~\cite{Radyushkin:2017cyf,Orginos:2017kos,Braun:2018brg,Li:2020xml}. The latter relies on
the validity of Euclidean operator product expansion (OPE) and can only be applied to
correlations at short distances, and therefore cannot be used directly for
LaMET applications.
In contrast, the RI/MOM method appears to be applicable to large-$z$ ($z$ is the gauge-link length)
distance at first glance. However, a detailed examination shows that this method introduces potential
non-perturbative effects, for instance, through infrared (IR) logarithms
such as $\ln(z^2\mu^2)$ ($\mu$ is a renormalization scale) in the scheme matching. Since UV divergences are supposed to be
perturbative in asymptotically-free theories such as QCD, it shall be possible to find a renormalization procedure which does not introduce  non-perturbative effects. It might be possible that RI/MOM does not introduce a large non-perturbative effect in the present precision
of lattice-QCD calculations. However, a systematic effective-theory
calculation with high precision cannot avoid addressing this issue.

To achieve this, we propose in this paper a hybrid renormalization procedure for lattice correlations in LaMET applications. At short distances
where OPE is valid,
the standard RI/MOM or ratio method is recommended.
At large distances, we suggest to use the auxiliary field formalism~\cite{Samuel:1978iy,Gervais:1979fv,Arefeva:1980zd,Dorn:1986dt} which has been advocated in LaMET
applications by a number of authors~\cite{Ji:2017oey,Green:2017xeu,Green:2020xco}.
In this formalism, the Wilson line is replaced by two-point functions of the auxiliary field. The linear divergence in
lattice correlation functions is then linked to the mass
renormalization of the auxiliary field, whereas the
logarithmic divergence appears in the renormalization of
the ``heavy''-light ``currents'' at the end of the Wilson line.
Both divergences can be separately renormalized in a manner which is consistent
with the $\overline{\rm MS}$ scheme~\cite{Green:2017xeu,
Green:2020xco}.  Although
the mass subtraction of the Wilson
line has been suggested before~\cite{Chen:2016fxx,Ishikawa:2017faj},
it has not been put into wide practical use because, to our knowledge,
a reliable approach to calculate the non-perturbative mass has not been well-established
in the literature. Here we suggest several ways to do so which shall be investigated through systematic
lattice simulations in the future.

In addition, we also address several other issues that are important in
extracting parton physics using LaMET, e.g., how to match appropriately to the continuum scheme near $z\sim 0$, and how to utilize the asymptotic behavior of relevant correlation functions at large light-front (LF) distance to remove the unphysical oscillations in the momentum distribution that arise from truncated Fourier transform.
